\chapter{NOTTI INQUIETE}\label{notti-inquiete}

\hfill\break

Seattle, settembre, cinque anni dopo il fallimento dell'attacco
missilistico cinese: guidavo verso casa nell'ora di punta, era venerdì e
pioveva; appena infilai la porta del mio appartamento, accesi
l'interfaccia audio e avviai una playlist che avevo creato e rinominato:
\emph{Terapia}.

Era stata una giornata lunga, al pronto soccorso dello Harborview. Mi
ero occupato di due casi di ferite da arma da fuoco e un tentativo di
suicidio. Sospesa dietro le palpebre mi era rimasta un'immagine di
sangue che colava dalle sponde di una barella. Mi tolsi i vestiti da
lavoro inumiditi di pioggia e indossai un paio di jeans e una felpa, mi
versai da bere e rimasi vicino alla finestra a guardare la città
sobbollire nel buio. Là fuori doveva esserci lo Stretto di Puget, un
vuoto senza luce oscurato da nuvole passeggere. Il traffico era quasi
fermo sull'Autostrada 5, un fiume rosso luminoso.

La mia vita, essenzialmente, era come me l'ero costruita. E stava in
equilibrio su una sola parola.

Di lì a poco, Astrud Gilberto cominciò a cantare con voce malinconica e
un po' stonata sulle note di \emph{Corcovado}, ma ero ancora troppo teso
per pensare a quello che mi aveva detto Jason al telefono la sera prima.
Troppo teso anche per ascoltare la musica nel modo in cui meritava di
essere ascoltata. \emph{Corcovado, Desafinado}, qualche brano di Gerry
Mulligan, qualcuno di Charlie Byrd. \emph{Terapia}. Ma nel rumore della
pioggia tutto perdeva i contorni. Scaldai la cena nel microonde e la
mangiai senza assaporarla; poi abbandonai ogni speranza di equilibrio
karmico e decisi di bussare alla porta di Giselle per vedere se era in
casa.

Giselle Palmer abitava in affitto a tre appartamenti dal mio. Aprì la
porta con indosso dei jeans strappati e una vecchia camicia di flanella
che preannunciavano una serata in casa. Le chiesi se aveva da fare o se
le andava di fare un giro.

«Non lo so, Tyler. Mi sembri piuttosto depresso.»

«Più che altro combattuto. Sto pensando di lasciare la città.»

«Veramente? Te ne vai, diciamo, in trasferta?»

«No, per sempre.»

«Ah» il suo sorriso svanì. «Quando l'hai deciso?»

«Non ho deciso. È questo il punto.»

Aprì di più la porta e mi fece segno di entrare. «Davvero? E dove
andresti?»

«È una lunga storia.»

«Vuol dire che devi bere qualcosa prima di parlarne?»

«Una cosa del genere.»

ooo

Avevo conosciuto Giselle l'anno prima, all'assemblea degli inquilini nel
seminterrato dell'edificio. Aveva ventiquattro anni e mi arrivava più o
meno all'altezza della spalla. Di giorno lavorava in un ristorante di
una qualche catena a Renton, ma quando iniziammo a trovarci per un caffè
la domenica pomeriggio mi spiegò che era ``una puttana, una prostituta,
è il mio lavoro part time''.

Intendeva dire che faceva parte di un gruppo disinvolto di amiche che si
scambiavano i nomi di uomini anziani (decorosi, solitamente sposati)
disposti a pagare generosamente per fare sesso ma terrorizzati dalla
prostituzione di strada. Mentre me lo raccontava, Giselle aveva drizzato
le spalle e mi guardava con aria di sfida, ipotizzando che potessi
esserne scioccato o ripugnato. Non avevo avuto nessuna di queste
reazioni. Dopo tutto, quelli erano gli anni dello Spin. Le persone
dell'età di Giselle creavano le proprie regole, nel bene e nel male, e
quelli come me si astenevano dal giudicare.

Continuammo a prendere il caffè o a cenare insieme di tanto in tanto, e
in un paio di occasioni le prescrissi la richiesta per le analisi del
sangue. Dall'ultimo esame, Giselle non era risultata sieropositiva e
l'unica malattia trasmissibile di cui aveva gli anticorpi era il virus
del Nilo occidentale. In altre parole, era stata attenta e fortunata.

Ma il punto era, mi aveva detto Giselle, che il sesso a pagamento, anche
a livello amatoriale, inizia a condizionarti la vita. «Diventi,» aveva
spiegato, «quel genere di persona che porta preservativi e Viagra nella
borsa.» Allora perché lo faceva, quando avrebbe potuto accettare, per
esempio, un lavoro notturno in un Walmart? Quella domanda non le era
piaciuta e aveva risposto stando sulla difensiva: «Forse è una
perversione. O forse è un hobby, sai, come quello dei trenini.» Ma
sapevo che era fuggita ancora bambina da Saskatoon, da un patrigno
violento, e la conseguente piega professionale non era difficile da
immaginare. E naturalmente, per i suoi comportamenti rischiosi, usava la
stessa scusa inoppugnabile che condividevamo tutti noi di una certa età:
la quasi certezza della nostra estinzione di massa. «La mortalità,»
disse una volta una scrittrice della mia generazione, «batte la
moralità.»

«Allora, a che livello alcolico vuoi arrivare? Solo brillo o
completamente sbronzo? In realtà, forse non abbiamo molta scelta.
Stasera il mobile bar è un po' sfornito.»

Mi fece un miscuglio, che era per lo più vodka e sembrava fosse stato
scolato da un serbatoio di benzina. Tolsi un quotidiano dalla sedia e mi
misi a sedere. L'appartamento di Giselle era arredato in modo decente ma
sembrava la stanza di una matricola in un dormitorio. Il giornale era
aperto alla pagina dell'editoriale. La vignetta umoristica riguardava lo
Spin: gli Ipotetici erano ritratti come una coppia di ragni neri che
stringono la Terra fra le zampe pelose. La didascalia recitava: {LI
	MANGIAMO SUBITO O ASPETTIAMO LE ELEZIONI}?

«Non ci capisco niente» disse Giselle, lasciandosi cadere sul divano e
indicando il giornale con un piede.

«Della vignetta?»

«Di tutto. Dello Spin. Di questo essere ``al punto di non ritorno''. Per
dire, mi spieghi a che serve leggere i giornali? C'è qualcosa dall'altra
parte del cielo, e non è nostro amico. Questo è quello che so e a me
basta.»

Probabilmente la maggior parte degli umani avrebbe apposto una firma
sotto quella dichiarazione. Ma per qualche motivo - forse per la pioggia
o il sangue versato quel giorno in mia presenza - le sue parole mi
indignarono. «Non è così difficile da capire.»

«No? Allora perché succede?»

«Il punto non è il \emph{perché}. Nessuno sa il \emph{perché}. Riguardo
al \emph{cosa}\ldots»

«Scusa, lo so, non c'è bisogno di farmi la lezione. Ci troviamo in una
specie di sacchetto cosmico e l'universo gira vorticosamente senza
controllo, bla bla bla.»

Le sue parole m'infastidirono di nuovo. «Conosci il tuo indirizzo, sì?»

Sorseggiò il suo drink. «Certo che lo conosco.»

«Perché ti piace sapere dove ti trovi. A circa tre chilometri
dall'oceano, a centocinquanta dal confine, a qualche migliaio a ovest di
New York, giusto?»

«Giusto, e quindi?»

«Cerco di fare un ragionamento. Le persone non hanno problemi a
distinguere Spokane da Parigi, ma quando si tratta del cielo vedono solo
una grande massa informe e misteriosa. Come mai?»

«Non lo so. Magari perché tutta l'astronomia che conosco l'ho imparata
dalle repliche di \emph{Star Trek}?. Voglio dire, ha davvero senso che
io sappia di lune e stelle? Cose che non vedo da quando ero bambina.
Perfino gli scienziati ammettono che spesso non sanno di che parlano.»

«E questo ti sta bene?»

«Che cazzo di differenza fa se sta bene a me? Ascolta, forse è meglio se
accendo la TV. Magari ci guardiamo un film e mi racconti perché stai
pensando di lasciare la città.»

«Le stelle sono come persone,» le dissi, «vivono e muoiono in intervalli
prevedibili. Il Sole sta invecchiando in fretta e brucia più velocemente
il suo combustibile. La sua luminosità aumenta del dieci per cento ogni
miliardo di anni. Il sistema solare è già cambiato tanto da rendere la
Terra naturale inabitabile, anche se lo Spin sparisse oggi stesso. Di
questo parlano i giornali quando scrivono ``punto di non ritorno''. E
non ci sarebbe stata nessuna notizia se il Presidente Clayton non
l'avesse reso ufficiale, ammettendo in un discorso che secondo i
migliori esperti scientifici non c'era modo di tornare allo \emph{status
	quo ante}.»

E lei mi rivolse un lungo sguardo infelice: «Tutte queste cazzate\ldots»

«Non sono \emph{cazzate}.»

«Forse no, ma non mi servono a niente.»

«Sto solo cercando di spiegare\ldots»

«Cazzo, Tyler. Ti ho chiesto una spiegazione? Prendi i tuoi incubi e
vattene a casa. Oppure rilassati e dimmi perché vuoi lasciare Seattle. È
per quegli amici tuoi, vero?»

Le avevo parlato di Jason e Diane. «Per Jason soprattutto.»

«Il cosiddetto genio.»

«Non solo cosiddetto. È in Florida\ldots»

«A fare qualcosa per le persone dei satelliti, hai detto.»

«A trasformare Marte in un giardino.»

«Era anche sui giornali. È davvero possibile?»

«Non ne ho idea. Jason pensa di sì.»

«Ma non ci vorrebbe molto tempo?»

«Gli orologi camminano più veloci dopo una certa altitudine.»

«A-ha, quindi per cosa ha bisogno di te?»

Be', sì, per cosa? Bella domanda. Ottima domanda.

«Alla Perielio devono assumere un medico per la clinica interna.»

«Pensavo fossi solo un medico generico.»

«Lo sono.»

«Quindi che qualifiche hai per essere un medico astronauta?»

«Assolutamente nessuna. Ma Jason\ldots»

«Sta facendo un favore a un vecchio amico? Be', allora torna tutto. Dio
benedica i ricchi, eh? Facciamo tutto in famiglia\ldots»

Scrollai le spalle. Che pensasse pure ciò che voleva. Non c'era bisogno
di condividere quelle cose con Giselle, e Jase non aveva specificato
niente.

Quando avevamo parlato, però, avevo avuto l'impressione che Jason non mi
volesse solo come medico dell'azienda ma anche come dottore personale.
Perché aveva un problema. Un genere di problema che non voleva
condividere con il personale della Perielio. Un problema di cui non
poteva parlare al telefono.

Giselle aveva finito la vodka ma rovistò nella borsa e tirò fuori uno
spinello nascosto in una scatola di assorbenti interni.

«La paga è buona, scommetto.» Sfregò la pietrina di un accendino di
plastica, avvicinò la fiamma alla punta attorcigliata dello spinello e
inspirò profondamente.

«Non siamo entrati nei dettagli.»

Espirò. «\emph{Quanto} sei geek. Forse è per questo che sopporti di
pensare continuamente allo Spin. Tyler Dupree, autistico borderline. È
\emph{questo} che sei, e lo sai. Ne hai tutti i sintomi. E scommetto che
questo Jason Lawton è proprio uguale a te. Scommetto che gli viene duro
ogni volta che pronuncia la parola ``miliardo''.»

«Non sottovalutarlo. Potrebbe davvero contribuire alla sopravvivenza
della specie umana.» O di qualche suo membro in particolare.

«La definirei un'ambizione da geek, se non fosse che non ne ho mai
sentita una. E questa sua sorella, quella con cui sei andato a
letto\ldots»

«Una volta.»

«Una volta. Era quella religiosa, giusto?»

«Giusto.» Lo era stata e lo era ancora, per quanto ne sapevo. Non avevo
notizie di Diane da quella notte nel Berkshire. E non perché non ci
avessi provato, almeno non solo per quello. A un paio di e-mail non
avevo ricevuto nessuna risposta. Neanche Jase l'aveva sentita molto ma
secondo Carol viveva con Simon in qualche zona dello Utah o dell'Arizona
- in qualche stato dell'Ovest che non avevo mai visitato e non riuscivo
a immaginarmi - dove erano rimasti bloccati dalla dissoluzione del
movimento del Nuovo Regno.

«Neanche questo è difficile da capire.» Giselle mi passò lo spinello.
Non ero proprio a mio agio con l'erba. E poi quell'osservazione sul
``geek'' mi aveva ferito. Feci un tiro profondo, e l'effetto fu lo
stesso che avevo già provato quando vivevo a Stony Brook: un'afasia
istantanea.

«Deve essere stato terribile per lei. L'evento dello Spin, lei che
voleva solo dimenticarsene, tu e la sua famiglia che non glielo
consentivate. Al posto suo anch'io sarei diventata di colpo religiosa.
Cazzo, a quest'ora anch'io starei cantando nel coro.»

«Il mondo è davvero così difficile da guardare?» biascicai sfasato e su
di giri.

Giselle allungò la mano e si riprese lo spinello. «Da dove sto io sì.
Quasi sempre.»

Girò la testa, distratta. Un tuono fece vibrare la finestra come se si
fosse risentito per il caldo secco dentro casa. Sullo stretto si
preparava un forte maltempo. «Scommetto che sarà uno di quegli inverni»
disse. «Di quelli brutti. Avrei voluto un camino. Magari la musica
sarebbe di aiuto. Ma sono troppo stanca per alzarmi.»

Mi avvicinai al suo impianto audio e avviai un album di Stan Getz che
aveva scaricato. Il sassofono riscaldò la stanza come nessun camino
avrebbe potuto fare. Lei fece un cenno con la testa: non era il disco
che avrebbe scelto lei, ma sì, andava bene.

«Quindi ti ha chiamato e ti ha offerto questo lavoro.»

«Giusto.»

«E gli hai detto che avresti accettato?»

«Gli ho detto che ci avrei pensato.»

«È quello che stai facendo? Ci stai pensando?»

Sembrava volesse insinuare qualcosa, ma non sapevo cosa. «Credo di sì.»

«Io credo di no. Credo che tu sappia già cosa farai. Sai cosa credo?
Credo che tu sia qui per dirmi addio.»

Le dissi che lo ritenevo possibile.

«Allora vieni almeno a sederti accanto a me.»

Mi spostai indolente sul divano. Giselle si stese e poggiò i piedi sulle
mie gambe. Indossava delle calze da uomo a rombi, lanuginose e un po'
ridicole. I risvolti dei jeans le salivano su per le caviglie. «Per
essere un tipo che riesce a guardare una ferita d'arma da fuoco senza
battere ciglio, sei abbastanza bravo a evitare gli specchi.»

«Che significa?»

«Significa che non hai chiuso davvero con Jason e Diane, è evidente.
Soprattutto con lei.»

Ma non era possibile che Diane fosse ancora importante per me.

Forse volevo dimostrarlo. Forse per questo ci ritrovammo nella camera
disordinata di Giselle, a fumare un altro spinello, a buttarci sul
copriletto rosa Barbie, a fare l'amore sotto le finestre oscurate dalla
pioggia, ad abbracciarci finché non ci addormentammo.

Ma non era il viso di Giselle che mi si presentò alla mente nella
vaghezza immediatamente successiva; mi svegliai un paio d'ore dopo
pensando: ``Oddio, ha ragione, andrò in Florida.''

ooo

Alla fine ci vollero settimane per organizzare tutto, sia per Jason che
per la clinica. In quel periodo, rividi Giselle ma solo di sfuggita.
Stava cercando una macchina usata da un rivenditore e così le vendetti
la mia; non volevo rischiare di attraversare il Paese in auto. (Il
numero delle rapine sulle interstatali era salito di molto.) Non facemmo
alcun cenno all'intimità che era venuta e se n'era andata con la
pioggia, un atto di gentilezza un po' brilla da parte di uno dei due,
più probabilmente sua.

A parte lei, a Seattle non avevo molte persone da salutare e neanche le
cose da portare via, escludendo dei file digitali, portatili per natura,
e qualche vecchio LP. Il giorno in cui partii, Giselle mi aiutò a
caricare i bagagli nel retro del taxi. «SeaTac» ordinai all'autista, e
mentre il taxi si infilava nel traffico lei mi salutò con la mano, non
particolarmente triste ma un po' malinconica.

Giselle era una brava persona che conduceva una vita pericolosa. Non la
rividi mai più ma spero che sia sopravvissuta al caos che sopraggiunse
in seguito.

ooo

Volai verso Orlando su un vecchio Airbus cigolante. Il rivestimento
della cabina era logoro e gli schermi sul dorso dei sedili ormai da
sostituire. Occupai il mio posto tra un uomo d'affari russo seduto
accanto al finestrino e una donna di mezza età sul lato corridoio. Il
russo, accigliato, si mostrò indifferente a qualsiasi conversazione, la
donna invece aveva voglia di chiacchierare: era una trascrittrice medica
diretta a Tampa per passare due settimane con la figlia e il genero. Si
chiamava Sarah, mi disse. Parlammo di dispositivi medici mentre l'aereo
si innalzava rumorosamente verso l'altitudine di crociera.

Nei cinque anni trascorsi dallo spettacolo dei fuochi d'artificio
cinesi, ingenti somme di denaro federale erano state immesse
nell'industria aerospaziale. Tuttavia, ne erano state dedicate
pochissime all'aviazione commerciale; ecco perché questi Airbus
risistemati volavano ancora. Il denaro era, invece, confluito in
progetti come quelli che E.D. Lawton gestiva dal suo ufficio di
Washington e Jason progettava alla Perielio in Florida: indagini sullo
Spin, fra cui la recente impresa su Marte. Era stata l'amministrazione
Clayton a convogliare quei finanziamenti, con la compiacenza di un
Congresso lieto di apparire impegnato a fare qualcosa di tangibile per
lo Spin. Ebbe effetti positivi sul morale pubblico. E meglio ancora,
nessuno si aspettava risultati concreti immediati.

Il denaro federale aveva contribuito a mantenere a galla l'economia
domestica, almeno nel Sud-Ovest, nell'area di Seattle, lungo le coste
della Florida. Ma era una prosperità infingarda e sottile come il
ghiaccio, e Sarah era preoccupata per sua figlia: suo genero era un
tubista autorizzato, sospeso a tempo indefinito dal lavoro presso un
distributore di gas naturale nella zona di Tampa. Vivevano in una
roulotte, prendevano un sussidio federale e cercavano di crescere un
bambino di tre anni, Buster, il nipote di Sarah.

«Non è un nome strano per un bambino? No, dico, \emph{Blister}? Sembra
una star del cinema muto. Ma il fatto è che gli si addice.»

Le dissi che i nomi sono come i vestiti: o sei tu a portare loro o sono
loro che portano te. «È proprio così, Tyler Dupree?» chiese, e io le
sorrisi imbarazzato.

«Comunque, non capisco proprio perché di questi tempi i giovani vogliano
avere figli. So che può sembrare orribile; non ho nulla contro Buster,
ovviamente. Gli voglio un bene dell'anima e spero abbia una vita lunga e
felice. Ma non posso fare a meno di chiedermi: che probabilità ci sono?»

«A volte le persone hanno bisogno di un motivo per sperare» le dissi,
chiedendomi se quella banale verità fosse ciò che Giselle aveva cercato
di dirmi.

«Però è vero anche che molti giovani \emph{non hanno} figli, intendo non
li hanno per scelta, come atto di generosità. Dicono che la cosa
migliore che si possa fare per un bambino è risparmiargli la sofferenza
che ci aspetta.»

«Non so se qualcuno sappia cosa ci aspetta.»

«Intendo il punto di non ritorno e tutto\ldots»

«Che abbiamo superato. Eppure siamo qui. Per qualche ragione.»

Inarcò le sopracciglia. «Crede che ci siano delle \emph{ragioni}, dottor
Dupree?»

Chiacchierammo ancora un po', poi Sarah disse: «Dovrei cercare di
dormire un po'» infilando il mini cuscino della compagnia aerea tra il
collo e il poggiatesta. Fuori dal finestrino, in parte nascosto
dall'indifferente russo, il Sole era tramontato, il cielo era diventato
nerofumo; non c'era niente da vedere se non il riverbero della luce da
lettura, che attenuai e orientai verso le ginocchia.

Come un idiota avevo messo tutta la roba da leggere nel bagaglio
imbarcato. Ma c'era una rivista sbrindellata nella tasca di fronte a
Sarah, e mi allungai per afferrarla. La rivista, con una semplice
copertina bianca, si chiamava \emph{Passaggio}. Una pubblicazione
religiosa, probabilmente lasciata lì da un precedente passeggero.

La sfogliai pensando inevitabilmente a Diane. Negli anni trascorsi dopo
l'attacco fallito agli artefatti dello Spin, il movimento del Nuovo
Regno aveva perso la coerenza di un tempo. Gli stessi fondatori lo
avevano sconfessato e il suo spensierato comunismo sessuale si era
spento sotto la pressione delle malattie veneree e della cupidigia
umana. Nessuno in quei giorni si sarebbe definito semplicemente un
membro del Nuovo Regno, nemmeno quelli che appartenevano ai movimenti
più estremi e marginali della religiosità alla moda. C'erano gli
ettoriani, i preteristi (completi o parziali), i ricostruzionisti del
Regno ma nessuno che fosse solo un membro dell'NR. Il circuito
dell'Ekstasis cui Diane e Simon avevano partecipato l'estate in cui ci
trovammo nel Berkshire non esisteva più.

Nessuna delle restanti fazioni NR aveva un peso demografico elevato. Già
solo i battisti del Sud superavano in numero tutte le sette del Regno
messe insieme. Ma il fulcro millenarista del movimento gli aveva
conferito un potere sproporzionato capace di alimentare l'ansia
religiosa che circondava lo Spin. Fu proprio a causa del Nuovo Regno che
molti tabelloni stradali annunciavano T{RIBOLAZIONE IN CORSO} e che
tante chiese tradizionali erano state costrette ad affrontare la
questione dell'apocalisse.

\emph{Passaggio} sembrava essere l'organo di stampa di una fazione
ricostruzionista della costa occidentale, destinata al grande pubblico.
Includeva, insieme a un editoriale che condannava calvinisti e
presbiteriani, tre pagine di ricette e una colonna di recensioni
cinematografiche. Ma quello che attirò la mia attenzione fu un articolo
intitolato I{L SACRIFICIO DI SANGUE E LA GIOVENCA ROSSA} - parlava di un
vitello rosso e innocente che sarebbe apparso ``compiendo la profezia''
e sacrificato al Monte del Tempio in Israele, annunciando l'Assunzione
in cielo. Evidentemente, l'antica fede NR nello Spin come atto di
redenzione non era più di moda. ``\emph{Come un laccio esso si abbatterà
	sopra tutti coloro che abitano sulla faccia di tutta la Terra}'', Luca
21,35. Un laccio, non una liberazione. Meglio trovare un animale da
sacrificare: la Tribolazione si stava rivelando più problematica del
previsto.

Rinfilai la rivista nella tasca del sedile mentre l'aereo incontrava
un'area di turbolenza. Sarah si accigliò nel sonno. L'uomo d'affari
russo suonò il campanello per chiamare lo steward e ordinare un whiskey
al succo di limone.

ooo

La macchina che noleggiai a Orlando il mattino seguente aveva due fori
di proiettile stuccati e verniciati ma ancora visibili sulla portiera
dal lato del passeggero. Chiesi al commesso se c'era qualcos'altro.

«È l'unica rimasta nel parcheggio, ma se non le dispiace attendere un
paio d'ore\ldots»

«No,» tagliai corto, «questa va bene.»

Imboccai l'Autostrada Bee Line verso est e poi svoltai verso sud sulla
95. Mi fermai a fare colazione lungo la strada in un Danny's fuori
Cocoa, dove la cameriera, forse intuendo la mia sostanziale condizione
di senzatetto, fu generosa con il caffè. «Viaggio lungo?»

«No, manca meno di un'ora.»

«Bene, praticamente sei arrivato. Vai a casa o te ne sei andato di
casa?»

Quando si rese conto che non avevo una risposta pronta mi sorrise.

«Sistemerai ogni cosa, caro. Lo facciamo tutti, prima o poi.» E in
cambio di quella benedizione da strada le lasciai una mancia fin troppo
generosa.

Il campus della Perielio - che Jason aveva chiamato, in maniera
allarmante, ``la roccaforte'' - era situato molto a sud delle
piattaforme di lancio di Canaveral-Kennedy, dove le sue strategie si
trasformavano in azioni fisiche. La Fondazione Perielio (ormai
ufficialmente un'agenzia governativa) non faceva parte della NASA,
sebbene si ``interfacciasse'' con essa, prendendo in prestito e
prestando ingegneri e personale. In un certo senso era uno strato di
burocrazia imposto alla NASA da tutte le amministrazioni successive allo
Spin, portando l'agenzia spaziale moribonda in direzioni che i suoi
vecchi capi non avrebbero potuto prevedere e forse non avrebbero
approvato. E.D. guidava il suo comitato direttivo e Jason aveva preso il
controllo effettivo sullo sviluppo del programma.

La giornata aveva cominciato a scaldarsi, un caldo tipico della Florida
che sembrava provenire dalla terra, la terra umida che trasudava come
una punta di petto in un barbecue. Passai accanto a distese di palme
nane frastagliate, negozi di surf sbiaditi, cunette piene di acqua verde
stagnante e una scena del crimine: auto della polizia che circondavano
un pick-up nero, tre uomini chini sul cofano di metallo rovente con i
polsi legati dietro la schiena. Un poliziotto che dirigeva il traffico
guardò a lungo la targa della mia auto e poi mi fece cenno di
proseguire, gli occhi gli brillavano di un sospetto vuoto e generico.

ooo

Quando raggiunsi la ``roccaforte'' Perielio, vidi che non era per niente
lugubre come poteva suggerire la parola. Era un complesso industriale
color salmone, moderno e pulito, immerso in un prato verde ondulato e
immacolato, i cui accessi erano ben vigilati ma poco intimidatori. Al
casotto una guardia sbirciò dentro la macchina, mi chiese di aprire il
bagagliaio, maneggiò in modo maldestro le valigie e la scatola con i
dischi, poi mi diede un pass temporaneo su una clip da taschino e mi
indirizzò verso il parcheggio dei visitatori (``dietro l'ala sud, segua
la strada alla sua sinistra, buona giornata''). La sua uniforme azzurra
era diventata indaco per il sudore.

Non feci in tempo a parcheggiare che vidi Jason attraversare un paio di
porte di vetro smerigliato contrassegnate da {TUTTI I VISITATORI DEVONO
	REGISTRARSI}, e un tratto di prato nel deserto rovente del parcheggio.
«Tyler!» esclamò, fermandosi bruscamente a un metro da me, come se
potessi svanire, quasi fossi un miraggio.

«Ehi, Jase» lo salutai sorridendo.

«Dottor Dupree!» ricambiò con un ampio sorriso. «Ma quella macchina? È a
noleggio? La faremo riportare a Orlando. Te ne daremo una più bella. Hai
già un posto dove stare?»

Gli ricordai che mi aveva garantito di occuparsi lui anche di questo.

«Ah, ce ne siamo occupati. O meglio, lo \emph{stiamo} facendo. Siamo in
trattativa per affittare un posticino a neanche venti minuti da qui.
Vista sull'oceano. Pronto in un paio di giorni. Nel frattempo avrai
bisogno di un hotel, ma a questo si può provvedere facilmente. Ma perché
dobbiamo starcene qui fuori ad assorbire raggi ultravioletti?»

Lo seguii nell'ala sud del complesso. Osservavo il modo in cui
camminava. Notai che pendeva un po' a sinistra e che preferiva usare la
mano destra.

Appena entrammo ci investì l'aria condizionata, un freddo artico che
dall'odore sembrava proveniente da cripte sterili e profonde. L'ingresso
era rivestito da numerose piastrelle lucide e granito. C'erano altre
guardie, queste sì addestrate a una gentilezza impeccabile. «Sono
proprio contento che tu sia qui» disse Jase. «Non dovrei perdere tempo,
ma ci tengo a farti vedere la struttura. Un giro rapido. In sala
conferenze mi aspettano delle persone della Boeing. Uno di Torrance e
uno del gruppo IDS di Saint Louis. Sono molto orgogliosi dei loro
potenziamenti di ioni allo xeno, che dovrebbero spremere un po' di
rendimento in più, come se fosse questa la cosa importante. Non ci
servono le sottigliezze, dico loro, ci servono affidabilità,
semplicità\ldots»

«Jason» lo interruppi.

«Loro\ldots{} che c'è?»

«Prendi fiato.»

Mi lanciò uno sguardo freddo e stizzito. Poi si addolcì, scoppiò a
ridere. «Scusami, è solo che, è come\ldots{} ricordi quando eravamo
bambini? Che ogni volta che uno di noi riceveva un giocattolo nuovo
doveva sfoggiarlo?»

Di solito era Jase quello che aveva i giocattoli nuovi, o almeno quelli
costosi. Comunque sì, gli dissi, mi ricordavo.

«Be', descriverlo in questo modo a chiunque altro sarebbe frivolo, però
qui, Tyler, abbiamo il più grande baule di giocattoli del mondo. Dai,
vieni che te lo mostro. Poi penseremo a sistemarti. Avrai il tempo di
adattarti al clima, ammesso che sia possibile.»

E così lo seguii per il pianterreno delle tre ali dell'edificio,
ammirando opportunamente le sale conferenze e gli uffici, gli enormi
laboratori e i reparti tecnici dove venivano concepiti i prototipi o
ridefiniti gli scopi delle missioni prima che la responsabilità di
progetti e obiettivi passasse ai grossi appaltatori. Tutto molto
interessante, tutto molto sconcertante. Infine arrivammo nell'infermeria
interna, dove mi presentò il dottor Koenig, il medico uscente, che mi
strinse la mano senza entusiasmo; poi, mentre andava via strascicando i
piedi, girò la testa per dirmi: «Buona fortuna, dottor Dupree.»

Nel frattempo il cercapersone tascabile di Jason continuava a suonare e
a quel punto non poté più ignorarlo. «Quelli della Boeing; mi tocca
ammirare le loro PPU, altrimenti mettono il broncio. Ce la fai a tornare
da solo alla reception? Lì ad aspettarti c'è Shelly, la mia assistente,
ti troverà lei una sistemazione. Possiamo continuare a parlare più
tardi. Tyler, è davvero bello rivederti!»

Un'altra stretta di mano, stranamente debole, poi si allontanò; vidi di
nuovo che pendeva a sinistra e non mi chiesi se fosse malato, ma quanto
lo fosse e quanto sarebbe potuto peggiorare.

ooo

Jason fu di parola. Nel giro di una settimana mi ero trasferito in una
piccola casa ammobiliata, apparentemente fragile come lo erano ai miei
occhi tutte le case della Florida, legno e assicelle, pareti fatte per
lo più di finestre, ma che doveva essere costata parecchio: il portico
rialzato si affacciava su una lunga discesa accanto a una via
commerciale che portava al mare. Durante quel periodo il taciturno
dottor Koenig mi ragguagliò ben tre volte sulle attività da svolgere;
era evidente che fosse stato infelice alla Perielio ma il passaggio di
consegne avvenne con grande serietà da parte sua. Mi affidò i suoi
fascicoli e il suo personale di supporto, e il lunedì successivo visitai
il mio primo paziente, un giovane metallurgo che si era storto la
caviglia durante una partita di calcio sul prato a sud. La clinica era
chiaramente ``sovrattrezzata'', come avrebbe detto Jase, per il lavoro
banale che facevamo ogni giorno. Ma Jason sosteneva che era meglio
muoversi in anticipo, perché forse un giorno sarebbe diventato difficile
ottenere cure mediche nel mondo oltre i cancelli.

Cominciai a sistemarmi. Scrivevo o prolungavo le prescrizioni,
distribuivo farmaci, sfogliavo i fascicoli. Scambiavo i soliti
convenevoli con Molly Seagram, la mia assistente, a cui ero risultato
molto più simpatico (disse lei) del dottor Koenig.

La sera me ne tornavo a casa e osservavo i fulmini guizzare dalle nuvole
attraccate al largo della costa come enormi velieri elettrizzati.

E aspettavo che Jason chiamasse: cosa che non fece mai, per quasi un
mese. Poi, un venerdì sera dopo il tramonto, si presentò
inaspettatamente alla porta, senza preavviso, in una tenuta sportiva
(jeans, maglietta) che gli toglieva dieci anni rispetto all'età che
dimostrava di solito. «Ho pensato di fare un salto, se per te non è un
problema.»

Certo che non lo era. Salimmo di sopra, andai a prendere due bottiglie
di birra dal frigorifero e ci sedemmo per un po' sul balcone imbiancato.
Jase cominciò a dirmi frasi del tipo «È fantastico rivederti» e «Che
bello averti a bordo» finché non lo interruppi: «Cazzo, Jase, non ho più
bisogno del comitato d'accoglienza. Sono semplicemente io.»

Rise imbarazzato, e poi tutto andò meglio.

Ci abbandonammo ai ricordi. A un certo punto gli chiesi: «La senti
spesso Diane?»

Si strinse nelle spalle. «Raramente.»

Non approfondii. Più tardi, quando ormai avevamo fatto fuori un paio di
birre a testa, l'aria era più fresca e la sera era diventata silenziosa,
gli chiesi come stava, intendendo a livello personale.

«Sono stato occupato, come avrai intuito. Siamo vicini ai primi lanci di
semi - più vicini di quanto abbiamo fatto credere alla stampa. A E.D.
piace giocare d'anticipo. Sta quasi sempre a Washington, Clayton in
persona ci tiene d'occhio. Siamo i prediletti dell'amministrazione,
almeno per ora. Ma questo mi costringe a occuparmi di questa merda
manageriale, che è infinita, e non del lavoro che voglio e \emph{devo}
fare, la progettazione della missione. È\ldots» con le mani fece un
gesto di impotenza.

«Stressante» completai la frase.

«Stressante. Ma stiamo facendo progressi. Passo dopo passo.»

«Ho notato che non c'è un fascicolo su di te. Alla clinica. Ho il
dossier sanitario di tutti gli altri dipendenti e amministratori, ma non
il tuo.»

Distolse lo sguardo, poi rise forte, una risata nervosa, secca.
«Be'\ldots{} mi piacerebbe lasciare le cose così, Tyler. Per il
momento.»

«Il dottor Koenig la pensava diversamente?»

«Il dottor Koenig pensa che siamo tutti un po' matti. E, ovviamente, è
vero. Te l'ho detto che ha accettato di gestire la clinica di una nave
da crociera? Te lo immagini? Koenig con la camicia hawaiana che
distribuisce antistaminici ai turisti?»

«Dimmi cosa c'è che non va, Jase.»

Fissò lo sguardo nel cielo orientale che si andava oscurando. C'era una
luce flebile sospesa qualche grado sopra l'orizzonte, non una stella,
quasi certamente uno degli aerostati di suo padre.

«Il fatto è,» disse quasi bisbigliando, «che ho un po' paura di essere
tagliato fuori proprio ora che iniziamo a ottenere dei risultati.» Mi
guardò a lungo. «Voglio potermi fidare di te, Ty.»

«Ci siamo solo noi, qui.»

Allora, finalmente, elencò i suoi sintomi - con calma, quasi
schematicamente, come se dolore e debolezza non fossero emotivamente più
importanti del grippaggio di un motore malfunzionante.

Gli promisi che gli avrei fatto fare degli esami ma non li avrei
inseriti tra le cartelle cliniche. Acconsentì annuendo, quindi lasciammo
perdere l'argomento e aprimmo ancora un'altra birra. Infine mi ringraziò
e mi strinse la mano, forse in maniera più solenne del necessario, e
lasciò la casa che aveva affittato per me, la nuova casa che sentivo
estranea.

Andai a letto temendo per lui.